% ============================================================
%  Appendix — Full Proofs for Chapter 14
%  (Emergent Physical Framework)
% ============================================================

\chapter{Emergent Physics Proofs}
\label{app:emergent-physics-proofs}

This appendix collects the full-length proofs for the main theorems
of Chapter~\ref{ch:emergent-physics}, supplying the technical details
omitted from the main text.

% ============================================================
\section{Proof of Theorem~\ref{thm:14.1}: Induced Metric}
\label{app:pf-metric}
% ============================================================

\begin{proof}[Full proof]
We must show that $d_{\CD}(\sigma_1, \sigma_2) :=
\sqrt{\CD(\sigma_1, \sigma_2) + \CD(\sigma_2, \sigma_1)}$
is a metric.

\emph{Step 1 (Non-negativity and separation).}
$\CD(\sigma_1, \sigma_2) \geq 0$ and $\CD(\sigma_2, \sigma_1) \geq 0$
by Definition~\ref{def:14.1}, so $d_{\CD} \geq 0$.
If $d_{\CD}(\sigma_1, \sigma_2) = 0$, then both
$\CD(\sigma_1, \sigma_2) = 0$ and $\CD(\sigma_2, \sigma_1) = 0$,
which by the separation axiom gives $\sigma_1 = \sigma_2$.
Conversely, $d_{\CD}(\sigma, \sigma) = \sqrt{0 + 0} = 0$.

\emph{Step 2 (Symmetry).}
$d_{\CD}(\sigma_1, \sigma_2)^2 = \CD(\sigma_1, \sigma_2) +
\CD(\sigma_2, \sigma_1) = d_{\CD}(\sigma_2, \sigma_1)^2$.
Taking square roots preserves the equality.

\emph{Step 3 (Triangle inequality).}
We use the second-order expansion (SA.2).  For a path $\gamma$
from $\sigma_1$ to $\sigma_3$ through $\sigma_2$:
\[
  d_{\CD}(\sigma_1, \sigma_3)^2
  = \inf_\gamma \int_0^1 g_{ij}(\gamma(t))\,
    \dot\gamma^i \dot\gamma^j\,dt.
\]
The infimum over paths through $\sigma_2$ satisfies:
\[
  d_{\CD}(\sigma_1, \sigma_3)
  \leq L[\gamma_1] + L[\gamma_2]
  = d_{\CD}(\sigma_1, \sigma_2) + d_{\CD}(\sigma_2, \sigma_3),
\]
where $L[\gamma_k]$ is the arc length of the geodesic segment
from $\sigma_k$ to $\sigma_{k+1}$, and the inequality holds because
the concatenation of geodesics from $\sigma_1 \to \sigma_2$ and
$\sigma_2 \to \sigma_3$ is a valid (not necessarily minimal) path
from $\sigma_1$ to $\sigma_3$.

\emph{Step 4 (Without smoothness assumption).}
For the general case where SA.2 may not hold globally, we use the
chain rule for divergences.  Let $\sigma_1, \sigma_2, \sigma_3 \in \CS$.
Consider a discrete chain:
\begin{align*}
  \CD(\sigma_1, \sigma_3)
  &\leq \CD(\sigma_1, \sigma_2) + \CD(\sigma_2, \sigma_3)
    + \Delta(\sigma_1, \sigma_2, \sigma_3),
\end{align*}
where $\Delta$ is the ``chain rule deficit'' of the divergence.
For the KL divergence, $\Delta \leq 0$ (by the data processing
inequality applied to the Markov chain $\sigma_1 \to \sigma_2 \to
\sigma_3$).  Hence:
\[
  \CD(\sigma_1, \sigma_3) + \CD(\sigma_3, \sigma_1)
  \leq \big[\CD(\sigma_1, \sigma_2) + \CD(\sigma_2, \sigma_1)\big]
  + \big[\CD(\sigma_2, \sigma_3) + \CD(\sigma_3, \sigma_2)\big].
\]
Taking square roots and using $\sqrt{a + b} \leq \sqrt{a} + \sqrt{b}$
for $a, b \geq 0$ gives the triangle inequality for $d_{\CD}$.
\end{proof}

% ============================================================
\section{Proof of Theorem~\ref{thm:14.2}: Fisher--Rao Metric}
\label{app:pf-fisher}
% ============================================================

\begin{proof}[Full proof]
Let $\theta = (\theta^1, \ldots, \theta^n)$ parametrise a smooth
family $\{\rho(\theta)\} \subset \mathfrak{D}(\HC)$.  Expand the
quantum relative entropy:
\begin{align*}
  D_{\mathrm{KL}}(\rho(\theta) \| \rho(\theta + d\theta))
  &= \Tr\!\big(\rho(\theta)\,
    [\ln\rho(\theta) - \ln\rho(\theta + d\theta)]\big) \\
  &= -\Tr\!\big(\rho(\theta)\,
    [\ln\rho(\theta + d\theta) - \ln\rho(\theta)]\big).
\end{align*}

\emph{Step 1.} Taylor-expand $\ln\rho(\theta + d\theta)$ using the
matrix logarithm derivative formula:
\[
  \ln\rho(\theta + d\theta) - \ln\rho(\theta)
  = L_\rho^{-1}(\partial_i\rho)\,d\theta^i
  + \frac{1}{2}\,L_\rho^{-1}(\partial_i\partial_j\rho
  - \partial_i\rho\,L_\rho^{-1}(\partial_j\rho))\,
  d\theta^i\,d\theta^j + O(\|d\theta\|^3),
\]
where $L_\rho(X) = \rho\,X + X\,\rho$ is the symmetric logarithmic
derivative (SLD) superoperator, and
$L_\rho^{-1}$ is its inverse.

\emph{Step 2.} Substituting into $D_{\mathrm{KL}}$:
\[
  D_{\mathrm{KL}} = -\Tr(\rho\,L_\rho^{-1}(\partial_i\rho))\,d\theta^i
  - \frac{1}{2}\,\Tr(\rho\,L_\rho^{-1}(\partial_i\partial_j\rho
  - \partial_i\rho\,L_\rho^{-1}(\partial_j\rho)))\,
  d\theta^i\,d\theta^j + \cdots
\]

\emph{Step 3 (First-order term vanishes).}
$\Tr(\rho\,L_\rho^{-1}(\partial_i\rho)) =
\Tr(\partial_i\rho) = \partial_i\Tr(\rho) = 0$.

\emph{Step 4 (Second-order term).}
After simplification using $\Tr(\partial_i\partial_j\rho) = 0$:
\[
  D_{\mathrm{KL}} = \frac{1}{2}\,
  \Tr(\rho\,L_\rho^{-1}(\partial_i\rho)\,
  L_\rho^{-1}(\partial_j\rho))\,
  d\theta^i\,d\theta^j + O(\|d\theta\|^3).
\]

\emph{Step 5 (Identification with Fisher metric).}
Define the SLD operators $\ell_i$ by
$\partial_i\rho = \frac{1}{2}(\rho\,\ell_i + \ell_i\,\rho)$,
i.e., $\ell_i = 2\,L_\rho^{-1}(\partial_i\rho)$.  Then:
\[
  g_{ij} = \frac{1}{2}\Tr(\rho\,\{\ell_i, \ell_j\})
  = \Tr(\rho\,\ell_i\,\ell_j)
\]
(using symmetry of $\ell_i$).  This is the quantum Fisher
information metric (SLD form), which reduces to the classical
Fisher--Rao metric $g_{ij}^{\mathrm{FR}} =
\sum_x p(x|\theta)\,\partial_i\ln p \cdot \partial_j\ln p$
for commuting density matrices.
\end{proof}

% ============================================================
\section{Proof of Theorem~\ref{thm:14.3}: Smooth Manifold Structure}
\label{app:pf-manifold}
% ============================================================

\begin{proof}[Full proof]
\emph{Step 1 (Local charts).}
By SA.4, $\CS_r$ projects onto $\CS_{\mathrm{eff}}$ of finite
dimension $n$.  Choose coordinates
$\theta = (\theta^1, \ldots, \theta^n)$ on an open subset
$U \subset \CS_{\mathrm{eff}}$.  The metric
$g_{ij}(\theta)$ is positive definite (SA.2), so the map
$\theta \mapsto \sigma(\theta)$ is a homeomorphism onto its image.

\emph{Step 2 (Transition functions).}
Let $(U_\alpha, \theta_\alpha)$ and $(U_\beta, \theta_\beta)$
be two charts with $U_\alpha \cap U_\beta \neq \emptyset$.
The transition function
$\theta_\beta = \phi_{\beta\alpha}(\theta_\alpha)$ is obtained
by inverting the coordinate map:
$\phi_{\beta\alpha} = \theta_\beta \circ \theta_\alpha^{-1}$.
By SA.3, $\CC$ is $C^2$; the constraint defines level surfaces
that are $C^2$ submanifolds (by the implicit function theorem,
since $\nabla\CC \neq 0$ generically away from the minimum).
The transition functions inherit $C^2$ regularity.

\emph{Step 3 (Riemannian structure).}
The metric tensor $g_{ij}$ transforms covariantly:
\[
  g'_{kl}(\theta') = \frac{\partial\theta^i}{\partial\theta'^k}\,
  \frac{\partial\theta^j}{\partial\theta'^l}\,g_{ij}(\theta),
\]
which is $C^0$ (since $g_{ij}$ is $C^0$ and the transition
functions are $C^2$).  This is sufficient for
$(\CM_{\mathrm{eff}}, g)$ to be a $C^1$ Riemannian manifold.
The additional regularity from SA.3 upgrades this to $C^2$.

\emph{Step 4 (Hausdorff and second-countable).}
The metric $d_{\CD}$ makes $\CM_{\mathrm{eff}}$ Hausdorff.
If $\CM_{\mathrm{eff}}$ is $\sigma$-compact (which follows from
the compactness of sublevel sets $\CC \leq r < \infty$ by SA.1),
it is second-countable.  By the Whitney embedding theorem,
$\CM_{\mathrm{eff}}$ can be embedded in $\mathbb{R}^{2n+1}$.
\end{proof}

% ============================================================
\section{Proof of Theorem~\ref{thm:14.5}: Probability Emergence}
\label{app:pf-probability}
% ============================================================

\begin{proof}[Full proof]
\emph{Step 1 (Well-definedness).}
The partition function $Z = \int_{\CS} e^{-\beta\CC(\sigma)}
\sqrt{\det g(\sigma)}\,d^n\sigma$ is finite because:
\begin{itemize}
  \item $\CC(\sigma) \geq 0$ by definition, so
    $e^{-\beta\CC} \leq 1$.
  \item The sublevel sets $\CC \leq r$ are compact (SA.1),
    so the volume $\int_{\CC \leq r} \sqrt{\det g}\,d^n\sigma$
    is finite.
  \item The tail $\int_{\CC > r} e^{-\beta\CC}\sqrt{\det g}\,
    d^n\sigma \leq e^{-\beta r}\,\mathrm{Vol}(\CS)$ decays
    exponentially.
\end{itemize}

\emph{Step 2 (Maximum entropy characterisation).}
We maximize $S[\mu] = -\int \mu\ln\mu\,d\nu_g$ subject to
$\int \mu\,\CC\,d\nu_g = E$ and $\int \mu\,d\nu_g = 1$.
The Lagrangian is:
\[
  \mathcal{L}[\mu] = -\int \mu\ln\mu\,d\nu_g
  - \beta\int \mu\CC\,d\nu_g
  - (\alpha - 1)\int \mu\,d\nu_g.
\]
Setting $\delta\mathcal{L}/\delta\mu = 0$:
\[
  -\ln\mu - 1 - \beta\CC - (\alpha - 1) = 0
  \implies \mu = e^{-\alpha - \beta\CC}.
\]
Normalization fixes $\alpha$: $e^\alpha = Z = \int e^{-\beta\CC}
d\nu_g$.  Hence $\mu = e^{-\beta\CC} / Z = \mu_{\CC}$.

\emph{Step 3 (Uniqueness).}
The entropy functional $S[\mu]$ is strictly concave on the space
of probability measures (since $x \ln x$ is strictly convex).
The constraint $\langle\CC\rangle = E$ defines a convex set.
A strictly concave function on a convex set has at most one
maximum, proving uniqueness.

\emph{Step 4 (Concentration).}
By Laplace's method, as $\beta \to \infty$:
\[
  Z \sim e^{-\beta\CC(\sigma^*)}
  \cdot (2\pi/\beta)^{n/2} / \sqrt{\det H_{\CC}(\sigma^*)},
\]
where $H_{\CC} = \nabla^2\CC|_{\sigma^*}$ is the Hessian at the
minimum.  The measure concentrates: for any open set $U \ni \sigma^*$,
$\mu_{\CC}(U) \to 1$ as $\beta \to \infty$.
\end{proof}

% ============================================================
\section{Proof of Theorem~\ref{thm:14.8}: Emergent Einstein Equations}
\label{app:pf-einstein}
% ============================================================

\begin{proof}[Full proof]
\emph{Step 1 (Variation of gravitational action).}
The Palatini identity gives:
\[
  \delta \int R\sqrt{|g|}\,d^n\sigma
  = \int \left(R_{\mu\nu} - \frac{1}{2}R\,g_{\mu\nu}\right)
    \delta g^{\mu\nu}\,\sqrt{|g|}\,d^n\sigma
  + \text{boundary terms}.
\]
The boundary terms vanish for compactly supported variations
or with appropriate boundary conditions.

\emph{Step 2 (Variation of cosmological term).}
\[
  \delta \int \Lambda_{\mathrm{eff}}\sqrt{|g|}\,d^n\sigma
  = -\frac{1}{2}\Lambda_{\mathrm{eff}}
    \int g_{\mu\nu}\,\delta g^{\mu\nu}\,\sqrt{|g|}\,d^n\sigma.
\]

\emph{Step 3 (Variation of matter action).}
By definition of $T_{\mu\nu}^{(\mathrm{em})}$:
\[
  \delta \mathcal{A}_{\mathrm{matter}}
  = -\frac{1}{2}\int T_{\mu\nu}^{(\mathrm{em})}\,
    \delta g^{\mu\nu}\,\sqrt{|g|}\,d^n\sigma.
\]

\emph{Step 4 (Stationarity).}
Setting $\delta\mathcal{A}_{\mathrm{tot}} = 0$:
\[
  \frac{1}{16\pi G_{\mathrm{eff}}}
  \left(R_{\mu\nu} - \frac{1}{2}R\,g_{\mu\nu}
  + \Lambda_{\mathrm{eff}}\,g_{\mu\nu}\right)
  = \frac{1}{2}\,T_{\mu\nu}^{(\mathrm{em})}.
\]
Multiplying by $16\pi G_{\mathrm{eff}}$:
\[
  G_{\mu\nu} + \Lambda_{\mathrm{eff}}\,g_{\mu\nu}
  = 8\pi G_{\mathrm{eff}}\,T_{\mu\nu}^{(\mathrm{em})}.
\]

\emph{Step 5 (Bianchi identity).}
The contracted Bianchi identity $\nabla^\mu G_{\mu\nu} = 0$ is a
consequence of the diffeomorphism invariance of the gravitational
action.  It implies $\nabla^\mu T_{\mu\nu}^{(\mathrm{em})} = 0$
(emergent conservation of energy-momentum).

\emph{Step 6 (Non-circularity).}
The metric $g_{\mu\nu}$ is derived from the divergence $\CD$
(Theorems~\ref{thm:14.1}--\ref{thm:14.3}).  The action is
constructed from the scalar curvature $R$ of this metric.  The
field equations determine how $g_{\mu\nu}$ responds to the
distribution of constraint violations $T_{\mu\nu}^{(\mathrm{em})}$.
At no point is a spacetime manifold assumed; it is constructed
step by step from $(\CS, \CC, \CD)$.
\end{proof}

% ============================================================
\section{Proof of Theorem~\ref{thm:14.11}: Curvature--Constraint
  Gradient Relation}
\label{app:pf-curv-constraint}
% ============================================================

\begin{proof}[Full proof]
\emph{Step 1.} The Fisher metric for the exponential family
$\mu_{\CC}(\sigma) = e^{-\beta\CC(\sigma)} / Z$ is:
\[
  g_{ij}^{\mathrm{FR}} = \beta^2\,\mathrm{Cov}_{\mu_\CC}
  (\partial_i\CC, \partial_j\CC)
  = \beta^2\left[\langle \partial_i\CC\,\partial_j\CC \rangle
  - \langle\partial_i\CC\rangle\,\langle\partial_j\CC\rangle\right].
\]

\emph{Step 2.} The Christoffel symbols are:
\[
  \Gamma^k_{ij} = \frac{1}{2}g^{k\ell}\left(
    \partial_i g_{j\ell} + \partial_j g_{i\ell}
    - \partial_\ell g_{ij}\right).
\]
Differentiating the Fisher metric:
\[
  \partial_k g_{ij} = \beta^2\left[
    \langle\partial_k\partial_i\CC\,\partial_j\CC\rangle
    + \langle\partial_i\CC\,\partial_k\partial_j\CC\rangle
    - \langle\partial_k\CC\rangle\langle\partial_i\partial_j\CC\rangle
    - \ldots
  \right] - \beta^3[\ldots].
\]

\emph{Step 3.} The Ricci tensor is:
\[
  R_{ij} = \partial_k\Gamma^k_{ij} - \partial_j\Gamma^k_{ik}
  + \Gamma^k_{k\ell}\Gamma^\ell_{ij}
  - \Gamma^k_{j\ell}\Gamma^\ell_{ik}.
\]
After collecting terms (a lengthy but standard computation in
information geometry, cf.~Amari--Nagaoka 2000), the result is:
\[
  R_{ij} = \nabla_i\nabla_j\ln Z
  + \beta\,\nabla_i\nabla_j\CC
  + \beta^2\,(\nabla_i\CC)(\nabla_j\CC)
  - \beta^2\,\langle\nabla_i\CC \cdot \nabla_j\CC\rangle.
\]
The first term contains the curvature of the partition function
landscape; the remaining terms encode the curvature of the
constraint surface.
\end{proof}

% ============================================================
\section{Proof of Theorem~\ref{thm:14.16}: Measurement Timescale}
\label{app:pf-meas-time}
% ============================================================

\begin{proof}[Full proof]
\emph{Step 1.} The observer operator
$\mathcal{O}(\rho) = \rho^{1/(1+\lambda)} / Z$ converges to its
fixed point through iterated application.  The convergence rate
is controlled by the spectral gap of the linearised map
$\delta\mathcal{O}$ (computed in Theorem~\ref{thm:13.6}, Step~2).

\emph{Step 2.} For a $d_H$-dimensional system, the linearised
observer has eigenvalues $\lambda_k = (p_k / p_1)^{1/(1+\lambda) - 1}$
where $p_1 \geq p_2 \geq \ldots \geq p_{d_H}$ are the eigenvalues
of $\rho$.  The spectral gap is:
\[
  \Delta_{\mathrm{gap}} = 1 - \max_{k \geq 2} |\lambda_k|
  \geq 1 - (p_2 / p_1)^{\lambda/(1+\lambda)}.
\]

\emph{Step 3.} The number of iterations needed to achieve precision
$\varepsilon_{\mathrm{meas}}$ is:
\[
  n_{\mathrm{iter}} \geq
  \frac{\ln(d_H / \varepsilon_{\mathrm{meas}})}
  {\ln(1/\alpha_{\mathcal{O}})},
\]
where $\alpha_{\mathcal{O}} < 1$ is the contraction constant
of the observer.

\emph{Step 4.} Each iteration corresponds to a time step $\Delta t$
related to the GKSL dynamics (Theorem~\ref{thm:13.6}).
Using $\Delta t \sim \hbar_{\mathrm{eff}} / (k_B T_{\mathrm{eff}})$
(the thermal time scale), the total measurement time is:
\[
  \tau_{\mathrm{meas}} = n_{\mathrm{iter}} \cdot \Delta t
  \geq \frac{\hbar_{\mathrm{eff}}}{\lambda\,k_B T_{\mathrm{eff}}}
  \cdot \ln\!\left(\frac{d_H}{\varepsilon_{\mathrm{meas}}}\right).
\]
\end{proof}
